\subsection{Regularized Yukawa kernels and path-length masses}

\[
P_{\mathrm{str}}:=P_{\Seam,+}P_{\mathrm{sing}}P_\Theta,
\qquad
P_{\Seam,+}:=\frac{\mathbf 1+\iC}{2},
\qquad
P_\Theta:=\frac12(\mathbf 1+Q_\Theta),
\qquad
P_{\mathrm{sing}}:=\frac13(\mathbf 1+Z_c+Z_c^2).
\]

\begin{definition}[Structural transport kernel]
\label{def:structural-transport-kernel}
For each fermion sector define the structural positive transport kernel by
\[
K_f^{\mathrm{str}}:=\bigl(P_{\mathrm{str}} D_f^\dagger D_f P_{\mathrm{str}}\bigr)^{-1/2},
\]
whenever the inverse exists on $\mathrm{Ran}(P_{\mathrm{str}})$. The final admissible kernel is
the restriction of this operator to the gap-stable branch selected below. The transport pole
$\delta_{\mathrm{ph}}$ is fixed by the structural stationarity of the exact determinant functional
\[
\mathcal F_{\mathrm{tr}}(\delta)
:=
\sum_{y\in\{1,1/3,2/3\}}
\Bigl[-\log\det\!\bigl(D_{y,\delta}^\dagger D_{y,\delta}+\varepsilon^2\bigr)\Bigr],
\]
with
\[
D_{y,\delta}:=y\,\mathbf 1-\delta U_6.
\]
On the closed branch one has
\begin{equation}
\label{eq:transport-bridge-offset}
\partial_\delta \mathcal F_{\mathrm{tr}}(\delta_{\mathrm{ph}})=0,
\qquad
\delta_{\mathrm{ph}}\in\left(\frac13,\frac23\right).
\end{equation}
\end{definition}

\begin{theorem}[Algebraic closure of the transport pole]
\label{prop:closed-form-transport-pole}
\label{thm:algebraic-transport-pole}
\label{thm:transport-pole-algebraic}
Let
\[
\mathcal F_{\mathrm{tr}}(\delta)
:=
\sum_{y\in\{1,1/3,2/3\}}
\Bigl[-\log\det\!\bigl(D_{y,\delta}^\dagger D_{y,\delta}+\varepsilon^2\bigr)\Bigr],
\qquad
D_{y,\delta}:=y\,\mathbf 1-\delta U_6.
\]
Write
\[
x:=\delta^2,
\qquad
\tau:=\varepsilon^2,
\qquad
z:=x^3=\delta^6.
\]
Then
\[
\det\!\bigl(D_{y,\delta}^\dagger D_{y,\delta}+\varepsilon^2\bigr)
=
Q_y(x,\tau),
\]
with
\[
Q_y(x,\tau)
=
\Bigl(x^2+2(\tau-y^2)x+(y^2+\tau)^2\Bigr)
\Bigl(x^2+(y^2+2\tau)x+(y^2+\tau)^2\Bigr)^2.
\]
Hence
\[
\partial_\delta \mathcal F_{\mathrm{tr}}(\delta)=0
\iff
\sum_{y\in\{1,1/3,2/3\}}
\frac{\partial_x Q_y(x,\tau)}{Q_y(x,\tau)}=0.
\]
On the closed zero-parameter branch one takes the intrinsic limit
\[
\tau=0
\]
after admissibility restriction. Then
\[
Q_y(x,0)=(y^6-z)^2,
\]
and the stationarity equation reduces to
\[
\frac{1}{1-z}
+
\frac{1}{1/729-z}
+
\frac{1}{64/729-z}
=0.
\]
Equivalently, for
\[
P(z):=
(z-1)\left(z-\frac{64}{729}\right)\left(z-\frac{1}{729}\right),
\]
one has
\[
P'(z)=0.
\]
The admissible transport pole is therefore the lower critical point of this cubic:
\[
z_\star
=
\frac{794-7\sqrt{9961}}{2187},
\qquad
\delta_{\mathrm{ph}}=z_\star^{1/6}
=
\left(\frac{794-7\sqrt{9961}}{2187}\right)^{1/6}
\approx \num{0.593277280133224}.
\]
As an expanded control form, one may equivalently write
\[
1594323\,\delta^{12}-1157652\,\delta^6+47449=0.
\]
\end{theorem}

\begin{proof}
In the Fourier basis of the six-site shift $U_6$, the operator
\[
D_{y,\delta}^\dagger D_{y,\delta}+\varepsilon^2
\]
has eigenvalues
\[
\lambda_{y,n}(x,\tau)
=
y^2+x-2y\sqrt{x}\cos\!\left(\frac{n\pi}{3}\right)+\tau,
\qquad n=0,\dots,5.
\]
Multiplying the $n=0,3$ factors and the conjugate pairs $(1,5)$ and $(2,4)$ gives exactly the
displayed factorization $Q_y(x,\tau)$. Since $x=\delta^2$ and $\delta\in(1/3,2/3)$ on the
admissible branch, one has $\delta>0$, so
\[
\partial_\delta\log Q_y(x,\tau)
=
2\delta\,\partial_x\log Q_y(x,\tau),
\]
and the stationarity condition is equivalent to the displayed $x$-equation. Setting $\tau=0$
collapses the factors to
\[
Q_y(x,0)
=(y^2-x)^2(y^4+xy^2+x^2)^2
=(y^6-z)^2.
\]
Therefore
\[
\sum_{y\in\{1,1/3,2/3\}}\frac{\partial_x Q_y(x,0)}{Q_y(x,0)}
=
-6x^2\sum_{y\in\{1,\frac13,\frac23\}}\frac{1}{y^6-z},
\]
so the stationary point condition reduces to the displayed rational equation. Let
\[
P(z):=(z-1)\left(z-\frac{64}{729}\right)\left(z-\frac{1}{729}\right).
\]
Since
\[
\frac{P'(z)}{P(z)}
=
\frac{1}{z-1}
+
\frac{1}{z-\frac{64}{729}}
+
\frac{1}{z-\frac{1}{729}},
\]
the stationary equation is equivalent to $P'(z)=0$. A real cubic with three ordered real roots has
exactly one critical point in each interval between consecutive roots, so the lower critical point
lies in
\[
z\in\left(\left(\frac13\right)^6,\left(\frac23\right)^6\right).
\]
Solving the quadratic $P'(z)=0$ gives
\[
z_\star=\frac{794-7\sqrt{9961}}{2187}.
\]
Hence
\[
\delta_{\mathrm{ph}}=z_\star^{1/6}\in\left(\frac13,\frac23\right).
\]
Expanding $P'(z)=0$ after substituting $z=\delta^6$ yields the displayed degree-$12$ control form.
\end{proof}

\begin{lemma}[Regularized transport pole stability]
\label{lem:regularized-transport-pole-stability}
Define the regularized stationarity function
\[
G(x,\tau):=
\sum_{y\in\{1,1/3,2/3\}}
\frac{\partial_x Q_y(x,\tau)}{Q_y(x,\tau)}.
\]
Let
\[
x_\star:=\delta_{\mathrm{ph}}^2,
\qquad
z_\star:=x_\star^3=\delta_{\mathrm{ph}}^6.
\]
Then
\[
\partial_xG(x_\star,0)\neq 0.
\]
Consequently there exists $\tau_0>0$ and a unique smooth branch
\[
x(\tau)\qquad (\tau\in[0,\tau_0))
\]
such that
\[
G(x(\tau),\tau)=0,
\qquad
x(0)=x_\star.
\]
Therefore
\[
\delta(\tau):=\sqrt{x(\tau)}=\delta_{\mathrm{ph}}+O(\tau)
\qquad (\tau\to 0).
\]
\end{lemma}

\begin{proof}
By the proof of \Cref{thm:transport-pole-algebraic},
\[
G(x,0)=6x^2\,\frac{P'(x^3)}{P(x^3)}.
\]
At $x=x_\star$ one has $P'(z_\star)=0$, so differentiating gives
\[
\partial_xG(x_\star,0)
=
18x_\star^4\,\frac{P''(z_\star)}{P(z_\star)}.
\]
Here $x_\star>0$, one has $P(z_\star)\neq 0$ because $z_\star$ is a critical point rather than a
root, and $P''(z_\star)\neq 0$ because the critical points of a cubic with three distinct roots are
simple. Hence $\partial_xG(x_\star,0)\neq 0$. The implicit-function theorem therefore yields a
unique smooth branch $x(\tau)$ through $(x_\star,0)$, and taking square roots gives the final
expansion for $\delta(\tau)$.
\end{proof}

\begin{definition}[Canonical cycle holonomy representation]
\label{def:canonical-cycle-holonomy}
On the rigid four-punctured family section
\[
X_f^\circ\cong \mathbb P^1\setminus\mu_4,
\qquad
\mu_4:=\{z\in\mathbb C:\ z^4=1\},
\]
fix a distinguished cycle basis
\[
(C_1,C_2,C_T)
\]
compatible with the lifted $D_4$ seam symmetry. Let
\[
F:=H^1(X_f^\circ,\mathbb C)\cong\mathcal H_F\cong\mathbb C^3.
\]
The canonical family holonomy datum is the determinant-normalized unitary representation
\[
\rho_F:\pi_1(X_f^\circ)\to SU(3)_F
\]
determined by the square symmetry of the puncture configuration, unimodularity, and the
admissible cusp assignment. Its concrete monodromy realization is forced by
\Cref{thm:d4-monodromy-rigidity}.
\end{definition}

\begin{theorem}[$D_4$-equivariant monodromy rigidity and Riemann--Hilbert closure]
\label{thm:d4-monodromy-rigidity}
\label{prop:generator-family-holonomy}
\label{thm:d4-riemann-hilbert-rigidity}
Choose positively oriented small loops
\[
\gamma_1,\gamma_i,\gamma_{-1},\gamma_{-i}\in\pi_1(X_f^\circ)
\]
around the punctures $1,i,-1,-i$, with the standard relation
\[
\gamma_1\gamma_i\gamma_{-1}\gamma_{-i}=1.
\]
Relative to the distinguished cycle basis $(C_1,C_2,C_T)$, set
\[
\omega:=e^{2\pi i/3},
\qquad
P:=
\begin{pmatrix}
0&1&0\\
0&0&1\\
1&0&0
\end{pmatrix},
\qquad
D:=
\begin{pmatrix}
1&0&0\\
0&\omega&0\\
0&0&\omega^2
\end{pmatrix}.
\]
If the determinant line is trivial, the local puncture spectra are
\[
\{1,\omega,\omega^2\},
\]
and the puncture square carries its inherited $D_4$ symmetry, then the admissible monodromy
representation is forced, up to global $SU(3)_F$ conjugation, by
\[
\rho_F(\gamma_1)=D,
\qquad
\rho_F(\gamma_i)=PDP^{-1},
\qquad
\rho_F(\gamma_{-1})=P^2DP^{-2},
\]
\[
\rho_F(\gamma_{-i})
:=
\bigl(\rho_F(\gamma_1)\rho_F(\gamma_i)\rho_F(\gamma_{-1})\bigr)^{-1}.
\]
Equivalently, the canonical monodromy package is no longer a generator input but a rigidity output
of the puncture-square symmetry, unimodularity, and admissible cusp data.
\end{theorem}

\begin{proof}
Because the local puncture spectrum is
\[
\{1,\omega,\omega^2\},
\qquad
\omega=e^{2\pi i/3},
\]
and the determinant line is trivial, every local monodromy matrix lies in $SU(3)_F$ and is
diagonalizable with three distinct eigenvalues. After one global $SU(3)_F$ conjugation we may
therefore assume
\[
\rho_F(\gamma_1)=D.
\]

Let $r$ denote the quarter-turn symmetry of the puncture square. Its action on the punctures sends
\[
1\mapsto i\mapsto -1\mapsto -i\mapsto 1,
\]
hence on the fundamental group it cyclically permutes the four puncture loops:
\[
r\gamma_1r^{-1}=\gamma_i,
\qquad
r\gamma_ir^{-1}=\gamma_{-1},
\qquad
r\gamma_{-1}r^{-1}=\gamma_{-i}.
\]
Because the family datum is $D_4$-equivariant, the same quarter-turn is implemented on the fiber by
an element of $SU(3)_F$ carrying the ordered eigenspace packet of $D$ cyclically onto itself. Since
the three eigenspaces are one-dimensional, a second global conjugation fixes that implementing
matrix to be the cyclic permutation matrix
\[
P=
\begin{pmatrix}
0&1&0\\
0&0&1\\
1&0&0
\end{pmatrix}.
\]
Therefore
\[
\rho_F(\gamma_i)=PDP^{-1},
\qquad
\rho_F(\gamma_{-1})=P^2DP^{-2}.
\]

The puncture loops satisfy exactly one relation,
\[
\gamma_1\gamma_i\gamma_{-1}\gamma_{-i}=1,
\]
so the last generator is forced:
\[
\rho_F(\gamma_{-i})
=
\bigl(\rho_F(\gamma_1)\rho_F(\gamma_i)\rho_F(\gamma_{-1})\bigr)^{-1}.
\]
Thus the monodromy representation is uniquely determined by the stated local spectra, determinant
triviality, and $D_4$ symmetry, up to the initial global $SU(3)_F$ conjugation.
\end{proof}

\begin{theorem}[Rigid admissible family holonomy class]
\label{thm:rigid-family-local-system}
On the rigid four-punctured family section
\[
X_f^\circ\cong \mathbb P^1\setminus\mu_4,
\]
fix the determinant line, the local monodromy conjugacy classes at the four punctures, and the
$D_4$ symmetry inherited from the seam winding theorem. Then there exists a unique admissible
flat Hermitian family connection
\[
\nabla_F^\star
\]
on the flat Hermitian family bundle $L_F\to X_f^\circ$ with fiber
$F\cong H^1(X_f^\circ,\mathbb C)\cong\mathbb C^3$ whose monodromy representation is the explicit
holonomy datum $\rho_F$ of
\Cref{def:canonical-cycle-holonomy,thm:d4-monodromy-rigidity}. Hence its $SU(3)_F$
conjugacy class $[\nabla_F^\star]$ is unique.
\end{theorem}
